\section{Introduction}

Gauss's \textit{Theorema Egregium \index{Theorema Egregium}} (1827) states that the Gaussian curvature \index{Gaussian curvature} of a surface is an intrinsic invariant: it can be determined entirely from measurements within the surface (the first fundamental form \index{fundamental form (differential geometry)}), without reference to the ambient space. This was a revolutionary discovery that founded differential geometry.

In this chapter, we cover:
\begin{itemize}
    \item First and second fundamental forms
    \item Gaussian curvature and mean curvature \index{mean curvature}
    \item The Theorema Egregium
    \item Geodesics and the Gauss-Bonnet \index{Gauss-Bonnet theorem} theorem
    \item Numerical computation of curvature
\end{itemize}

\section{Surfaces and Fundamental Forms}

Let $\mathbf{r}(u,v)$ be a parametrized surface \index{surface (differential geometry)} in $\R^3$.

\begin{definition}[First Fundamental Form (Metric Tensor)]
\[
I = E du^2 + 2F du dv + G dv^2,
\]
where
\[
E = \mathbf{r}_u \cdot \mathbf{r}_u,\quad
F = \mathbf{r}_u \cdot \mathbf{r}_v,\quad
G = \mathbf{r}_v \cdot \mathbf{r}_v.
\]
\end{definition}

\begin{definition}[Second Fundamental Form]
Let $\mathbf{N} = \frac{\mathbf{r}_u \times \mathbf{r}_v}{\|\mathbf{r}_u \times \mathbf{r}_v\|}$ be the unit normal \index{normal vector}.
\[
II = L du^2 + 2M du dv + N dv^2,
\]
where
\[
L = \mathbf{r}_{uu} \cdot \mathbf{N},\quad
M = \mathbf{r}_{uv} \cdot \mathbf{N},\quad
N = \mathbf{r}_{vv} \cdot \mathbf{N}.
\]
\end{definition}

\section{Curvature}

\begin{definition}[Gaussian and Mean Curvature]
\[
K = \frac{LN - M^2}{EG - F^2}, \quad
H = \frac{EN - 2FM + GL}{2(EG - F^2)}.
\]
\end{definition}

\begin{theorem}[Theorema Egregium]
The Gaussian curvature $K$ depends only on the first fundamental form $E, F, G$ and their derivatives. It is an intrinsic invariant of the surface.
\end{theorem}

\begin{proof}[Proof Sketch]
The Christoffel symbols \index{Christoffel symbol} $\Gamma_{ij}^k$ are computed from $E, F, G$. The Riemann curvature tensor \index{Riemann curvature tensor} $R_{ijkl}$ is expressed in terms of $\Gamma$. The Gaussian curvature is $K = R_{1212} / \det(g)$. No dependence on $L, M, N$ remains.
\end{proof}

\section{Christoffel Symbols and Riemann Tensor}

\begin{definition}[Christoffel Symbols \index{Christoffel symbol}]
The Christoffel symbols of the second kind are:
\[
\Gamma_{ij}^k = \frac{1}{2} g^{kl} \left(\frac{\partial g_{il}}{\partial x^j} + \frac{\partial g_{jl}}{\partial x^i} - \frac{\partial g_{ij}}{\partial x^l}\right).
\]
\end{definition}

For the first fundamental form $g_{11} = E$, $g_{12} = F$, $g_{22} = G$:
\begin{align*}
\Gamma_{11}^1 &= \frac{GE_u - 2FF_u + FE_v}{2(EG-F^2)} \\
\Gamma_{12}^1 &= \frac{GE_v - FG_u}{2(EG-F^2)} \\
\Gamma_{22}^1 &= \frac{2GF_v - GG_u - FG_v}{2(EG-F^2)} \\
\Gamma_{11}^2 &= \frac{2EF_u - EE_v - FE_u}{2(EG-F^2)} \\
\Gamma_{12}^2 &= \frac{EG_u - FE_v}{2(EG-F^2)} \\
\Gamma_{22}^2 &= \frac{EG_v - 2FF_v + FG_u}{2(EG-F^2)}
\end{align*}

\begin{definition}[Riemann Curvature Tensor \index{Riemann curvature tensor}]
The Riemann curvature tensor measures the failure of parallel transport to be path-independent:
\[
R^l_{ijk} = \frac{\partial \Gamma_{ik}^l}{\partial x^j} - \frac{\partial \Gamma_{ij}^l}{\partial x^k} + \Gamma_{im}^l \Gamma_{jk}^m - \Gamma_{km}^l \Gamma_{ij}^m.
\]
\end{definition}

Then $R_{1212} = \frac{1}{2}(E_{vv} + G_{uu} - 2F_{uv}) + \text{terms in } \Gamma$.

\section{Python Implementation}

In \texttt{python/theorema.py}:

\begin{lstlisting}[style=python, caption={Key functions in python/theorema.py}]
first_fundamental_form(r, u, v)     # Returns E, F, G
second_fundamental_form(r, u, v)    # Returns L, M, N
christoffel_symbols(E, F, G)        # Returns Gamma^k_ij
gaussian_curvature(E, F, G, ...)    # K from first form only
mean_curvature(E, F, G, L, M, N)    # H
riemann_tensor(E, F, G)             # R_ijkl
theorema_egregium_check(r, u, v)    # Verify K intrinsic
geodesic \index{geodesic}_equations(r)               # Geodesic ODEs
gauss_bonnet_discrete(vertices, faces)  # Discrete Gauss-Bonnet
\end{lstlisting}

\section{Numerical Examples}

\begin{example}[Sphere of Radius $R$]
$\mathbf{r}(\theta,\phi) = (R\sin\theta\cos\phi, R\sin\theta\sin\phi, R\cos\theta)$
$E = R^2$, $F = 0$, $G = R^2\sin^2\theta$, $K = 1/R^2$, $H = 1/R$.
\begin{lstlisting}[style=python]
>>> from theorema import gaussian_curvature, first_fundamental_form
>>> import numpy as np
>>> r = lambda u,v: np.array([np.sin(u)*np.cos(v), np.sin(u)*np.sin(v), np.cos(u)])
>>> E, F, G = first_fundamental_form(r, 1.0, 0.5)
>>> K = gaussian_curvature(r, 1.0, 0.5)
>>> print(f"K = {K:.6f}")  # 1.0
\end{lstlisting}
\end{example}

\begin{example}[Cylinder]
$\mathbf{r}(\theta,z) = (R\cos\theta, R\sin\theta, z)$
$K = 0$ (intrinsically flat), $H = 1/(2R)$.
\end{example}

\begin{example}[Pseudosphere]
Surface of revolution of the tractrix: $K = -1$ (constant negative curvature).
\end{example}

\section{Exercises}

\begin{exercise}
Prove that the Gaussian curvature of a surface of revolution $z = f(r)$ is $K = \frac{f' f''}{r(1+f'^2)^2}$.
\end{exercise}

\begin{exercise}
Implement discrete differential geometry operators on triangle meshes to compute $K$ at vertices.
\end{exercise}

\begin{exercise}
Verify the Theorema Egregium numerically: compute $K$ from first form only and compare with $K$ from both forms for various surfaces.
\end{exercise}

\begin{exercise}
Apply the discrete Gauss-Bonnet theorem to a triangulated torus and verify $\sum K_i A_i = 0$.
\end{exercise}

\begin{exercise}
Compute the geodesics on a sphere of radius $R$ and verify they are great circles.
\end{exercise}



\section{Exercise Solutions}

\begin{solution}
The Theorema Egregium states that Gaussian curvature $K = \kappa_1 \kappa_2$ is invariant under isometric deformations. This follows from the Gauss equation $R_{1212} = K(g_{11}g_{22} - g_{12}^2)$, where the Riemann curvature tensor depends only on the first fundamental form and its derivatives.
\end{solution}

\begin{solution}
For a surface of revolution $z = f(r)$, parametrize as $\mathbf{r}(r,\theta) = (r\cos\theta, r\sin\theta, f(r))$. Then $E = 1 + f'^2$, $F = 0$, $G = r^2$. Computing Christoffel symbols and the Riemann tensor gives $K = \frac{f' f''}{r(1+f'^2)^2}$.
\end{solution}

\begin{solution}
For a triangle mesh, the discrete Gaussian curvature at vertex $v_i$ is $K_i = \frac{2\pi - \sum_{j \in \text{stars}(i)} \theta_j}{A_i}$, where $\theta_j$ are the face angles at $v_i$ and $A_i$ is the circumscript area.
\end{solution}

\begin{solution}
For a triangulated torus, assign curvature $K_i = 2\pi - \sum \theta_j$ at each vertex (angle deficit). The discrete Gauss-Bonnet theorem states $\sum_i K_i A_i = 2\pi\chi(T^2) = 0$. For a rectangular torus mesh, each interior vertex has angle sum $2\pi$, so $K_i = 0$.
\end{solution}

\begin{solution}
Geodesics on a sphere satisfy the geodesic equation $\ddot{u}^k + \Gamma_{ij}^k \dot{u}^i \dot{u}^j = 0$. For a sphere, the solutions are great circles: $\theta(t) = at + b$, $\phi(t) = c$. The geodesic distance is $R \cdot \arccos(\sin\theta_1\sin\theta_2 + \cos\theta_1\cos\theta_2\cos(\phi_1-\phi_2))$.
\end{solution}


\section{References}

\begin{itemize}
    \item \cite{gauss-disquisitiones-generales} -- Gauss, \textit{Disquisitiones Generales circa Superficies Curvas} (presented 1827; publ. 1828)
    \item \cite{do-carmo} -- do Carmo, \textit{Differential Geometry of Curves and Surfaces}
    \item \cite{spivak-dg} -- Spivak, \textit{A Comprehensive Introduction to Differential Geometry}
    \item \cite{meyer-discrete} -- Meyer et al., \textit{Discrete Differential-Geometry Operators for Triangulated 2-Manifolds}
\end{itemize}